\subsection{Homogeneous polynomial mappings}

PROPOSITION 13. --- Let M and N be A-modules, q an integer ≥ 0, and f a mapping of M into N. Suppose that M is free, then the following conditions are equivalent:

\begin{enumerate}
\def\labelenumi{(\roman{enumi})}
\item
  There exists a \(q\) -linear mapping \(g\) of \(\mathbf{M}^q\) into \(\mathbf{N}\) such that \(f(x) = g(x, x, \ldots, x)\) for all \(x \in \mathbf{M}\) .
\item
  There exists a linear mapping \(h\) of \(\mathbf{T}\mathbf{S}^q (\mathbf{M})\) into \(\mathbf{N}\) such that \(f(x) = h(\gamma_q(x))\) for all \(x \in \mathbf{M}\) .
\item
  There exists a basis \((e_i)_{i\in I}\) of \(\mathbf{M}\) and a family \((u_{\nu})_{\nu \in \mathbf{N}^{(I)},|\nu| = q}\) of elements of \(\mathbf{N}\) such that
\end{enumerate}

\[
f \left(\sum_ {i \in I} \lambda_ {i} e _ {i}\right) = \sum_ {\nu \in \mathbf {N} ^ {(I)}, | \nu | = q} \lambda^ {\nu} u _ {\nu}\tag{15}
\]

for all \((\lambda_i) \in \mathbf{A}^{(I)}\) .

\begin{enumerate}
\def\labelenumi{(\roman{enumi})}
\setcounter{enumi}{3}
\tightlist
\item
  For each basis \((e_i)_{i\in I}\) of \(\mathbf{M}\) there exists a family \((u_{\nu})_{\nu\in \mathbf{N}^{(I)},|\nu| = q}\) of elements of \(\mathbf{N}\) such that
\end{enumerate}

\[
f \left(\sum_ {i \in I} \lambda_ {i} e _ {i}\right) = \sum_ {\nu \in \mathbf {N} ^ {(I)}, | \nu | = q} \lambda^ {\nu} u _ {\nu}
\]

for all \((\lambda_i) \in \mathbf{A}^{(\mathrm{I})}\) .

\begin{enumerate}
\def\labelenumi{(\roman{enumi})}
\tightlist
\item
  \(\Rightarrow\) (ii): let \(g\) satisfy (i), then there exists a linear mapping \(g'\) of \(\mathbf{T}^q (\mathbf{M})\) into \(\mathbf{N}\) such that \(g(x_{1},x_{2},\ldots ,x_{q}) = g'(x_{1}\otimes x_{2}\otimes \ldots \otimes x_{q})\) for any \(x_{1},\ldots ,x_{q}\in \mathbf{M}\) . Then
\end{enumerate}

\[
f (x) = g (x, x, \dots , x) = g ^ {\prime} (x \otimes x \otimes \dots \otimes x) = g ^ {\prime} (\gamma_ {q} (x));
\]

and on writing \(h = g' \mid \mathbf{TS}^q(\mathbf{M})\) we see that condition (ii) holds.

\begin{enumerate}
\def\labelenumi{(\roman{enumi})}
\setcounter{enumi}{1}
\tightlist
\item
  \(\Rightarrow\) (i) and (iv): let \(h\) satisfy the conditions of (ii). By Prop. 4, (ii) (IV, p.~47) there exists a linear mapping \(g'\) of \(\mathbf{T}^q (\mathbf{M})\) into \(\mathbf{N}\) such that \(h = g'|\mathbf{TS}^q (\mathbf{M})\) . Let \(g\) be the \(q\) -linear mapping of \(\mathbf{M}\) into \(\mathbf{N}\) associated with \(g'\) , then for any \(x\in \mathbf{M}\) we have
\end{enumerate}

\[
f (x) = h \left(\gamma_ {q} (x)\right) = g ^ {\prime} (x \otimes x \otimes \dots \otimes x) = g (x, x, \dots , x),
\]

whence (i). On the other hand, if \((e_i)_{i\in I}\) is a basis of \(\mathbf{M}\) , we have, by formula (6) (IV, p.~46)

\[
f \left(\sum_ {i} \lambda_ {i} e _ {i}\right) = h \left(\gamma_ {q} \left(\sum_ {i} \lambda_ {i} e _ {i}\right)\right) = h \left(\sum_ {| \nu | = q} \lambda^ {\nu} e _ {\nu}\right)
\]

on writing \(e_{\nu} = \prod_{i\in I}\gamma_{\nu_i}(e_i)\) ; we thus have

\[
f \left(\sum_ {i} \lambda_ {i} e _ {i}\right) = \sum_ {| \nu | = q} \lambda^ {\nu} h \left(e _ {\nu}\right).
\]

\begin{enumerate}
\def\labelenumi{(\roman{enumi})}
\setcounter{enumi}{3}
\item
  \(\Rightarrow\) (iii) is clear.
\item
  \(\Rightarrow\) (ii): given \((e_i)\) , \((u_\nu)\) satisfying the conditions of (iii), let us write \(e_\nu = \prod_{i\in I}\gamma_{\nu_i}(e_i)\) and recall that \((e_\nu)_{|\nu| = q}\) is a basis of \(\mathbf{TS}^q (\mathbf{M})\) . Let \(h\) be the homomorphism of \(\mathbf{TS}^q (\mathbf{M})\) into \(\mathbf{N}\) defined by \(h(e_{\nu}) = u_{\nu}\) ; then for each \(x = \sum_{i\in I}\lambda_i e_i\) in \(\mathbf{M}\) we have
\end{enumerate}

\[
f (x) = f \left(\sum_ {i} \lambda_ {i} e _ {i}\right) = \sum_ {| \nu | = q} \lambda^ {\nu} u _ {\nu} = h \left(\sum_ {| \nu | = q} \lambda^ {\nu} e _ {\nu}\right) = h (\gamma_ {q} (x)).
\]

DEFINITION 3. --- Let M and N be A-modules and q an integer \(\geqslant 0\) . Suppose that M is free, and denote by \(\operatorname{Pol}_{\mathbf{A}}^{q}(\mathbf{M}, \mathbf{N})\) or simply \(\operatorname{Pol}^{q}(\mathbf{M}, \mathbf{N})\) the set of mappings of M into N satisfying the conditions of Prop. 13. The elements of \(\operatorname{Pol}^{q}(\mathbf{M}, \mathbf{N})\) are called homogeneous polynomial mappings of degree q of M into N.

Prop. 13 (i) defines a homomorphism of A-modules :

\[
\mathcal {L} _ {q} (\mathrm{M}, \dots , \mathrm{M}; \mathrm{N}) \rightarrow \operatorname{Pol} ^ {q} (\mathrm{M}, \mathrm{N}).
\]

Prop. 13 (ii) defines a homomorphism of A-modules :

\[
\operatorname{Hom} _ {\mathrm{A}} \left(\mathbf {T S} ^ {q} (\mathrm{M}), \mathrm{N}\right)\rightarrow \operatorname{Pol} ^ {q} (\mathrm{M}, \mathrm{N}).
\]

These homomorphisms are called canonical. They are surjective.

Examples. --- 1) The homogeneous polynomial mappings of degree 1 of M into N are the A-linear mappings of M into N.

\begin{enumerate}
\def\labelenumi{\arabic{enumi})}
\setcounter{enumi}{1}
\tightlist
\item
  Let \((\mathbf{N}_i)_{i\in I}\) be a family of A-modules, \(f_{i}\) a mapping of \(\mathbf{M}\) into \(\mathbf{N}_i\) , \(i\in I\) , and \(f:\mathbf{M}\to \prod_{i\in I}\mathbf{N}_i\) the mapping with components \(f_{i}\) . For \(f\) to be a
\end{enumerate}

homogeneous polynomial mapping of degree q it is necessary and sufficient that each \(f_{i}\) be a homogeneous polynomial mapping of degree q.

\begin{enumerate}
\def\labelenumi{\arabic{enumi})}
\setcounter{enumi}{2}
\item
  Let \((\mathbf{M}_j)_{j\in J}\) be a finite family of free A-modules and \(u:\prod_{j\in J}\mathbf{M}_j\to \mathbf{N}\) a multilinear mapping; then \(u\) is polynomial of degree Card(J).
\item
  Let \((\mathbf{X}_i)_{i\in I}\) be a family of indeterminates, N an A-module and \(u\in \mathbf{N}[(\mathbf{X}_i)_{i\in I}]\) a homogeneous polynomial of degree \(q\) . The mapping \((x_{i})_{i\in I}\mapsto u((x_{i})_{i\in I})\) of \(\mathbf{A}^{(I)}\) into N is a homogeneous polynomial mapping of degree \(q\) : this is seen at once by condition (iii) of Prop. 13. If I is finite, every homogeneous polynomial mapping of degree \(q\) of \(\mathbf{A}^{(I)} = \mathbf{A}^{\mathrm{I}}\) into N is of that form.
\item
  The mapping \((x_{i})_{i\in\mathbf{N}}\mapsto x_{0}^{2}+x_{1}^{2}+\cdots+x_{n}^{2}+\cdots\) of \(\mathbf{A}^{(N)}\) into A is a homogeneous polynomial mapping of degree 2. If A=Z/2Z, it coincides with the linear mapping \((x_{i})_{i\in I}\mapsto x_{0}+x_{1}+\cdots+x_{n}+\cdots\)
\item
  Let \(f \in \operatorname{Pol}_{\mathbf{A}}^{q}(\mathbf{M}, \mathbf{N})\) , let \(\mathbf{B}\) be a commutative ring, \(\rho\) a homomorphism of \(\mathbf{B}\) into \(\mathbf{A}\) and \(\mathbf{M}'\) and \(\mathbf{N}'\) the B-modules derived from \(\mathbf{M}\) and \(\mathbf{N}\) by means of \(\rho\) . Assume that \(\mathbf{M}'\) is free; then \(f \in \operatorname{Pol}_{\mathbf{B}}^{q}(\mathbf{M}', \mathbf{N}')\) : this follows at once from condition (i) of Prop. 13.
\end{enumerate}

PROPOSITION 14. --- Let M, N, P be A-modules, q and r integers ≥ 0, and assume that M and N are free. If \(f \in \text{Pol}^{q}(M, N)\) , \(f' \in \text{Pol}^{r}(N, P)\) , then \(f' \circ f \in \text{Pol}^{qr}(M, P)\) .

There exists a q-linear mapping g of \(M^{q}\) into N and an r-linear mapping \(g'\) of \(N^{r}\) into P such that

\[
\begin{array}{l l} f (x) = g (x, x, \dots , x) & \text { for   all } x \in \mathbf {M}, \\ f ^ {\prime} (y) = g ^ {\prime} (y, y, \dots , y) & \text { for   all } y \in \mathbf {N}. \end{array}
\]

Hence for every \(x \in \mathbf{M}\) we have

\[
f ^ {\prime} (f (x)) = g ^ {\prime} (f (x), f (x), \dots , f (x)) = g ^ {\prime} (g (x, x, \dots , x), \dots , g (x, x, \dots , x))
\]

and the mapping \((x_{1},\ldots ,x_{qr})\mapsto g^{\prime}\big(g(x_{1},\dots,x_{q}),\dots,g(x_{q(r - 1) + 1},\dots,x_{qr})\big)\) of \(\mathbf{M}^{qr}\) into \(\mathbf{P}\) is \(qr\) -linear.

PROPOSITION 15. --- Let M be a free A-module, N an A-module and q an integer \(\geqslant0\) . We suppose that the mapping \(y\mapsto q!y\) is an automorphism of N. Let \(f\in\mathrm{Pol}^{q}(\mathbf{M},\mathbf{N})\) , then there exists one and only one symmetric q-linear mapping h of \(M^{q}\) into N such that \(f(x)=h(x,x,\ldots,x)\) for all \(x\in M\) . For any \(x_{1},\ldots,x_{q}\in M\) we have

\[
h \left(x _ {1}, x _ {2}, \dots , x _ {q}\right) = \frac {(- 1) ^ {q}}{q !} \sum_ {\mathrm{H} \subset \{1, 2, \dots , q \}} (- 1) ^ {\text { Card   H }} f \left(\sum_ {i \in \mathrm{H}} x _ {i}\right).\tag{16}
\]

\begin{enumerate}
\def\labelenumi{\alph{enumi})}
\tightlist
\item
  There exists a \(q\) -linear mapping \(g\) of \(\mathbf{M}^q\) into \(\mathbf{N}\) such that \(f(x) = g(x, x, \ldots, x)\) for all \(x \in \mathbf{M}\) . Let us define a \(q\) -linear mapping \(h\) of \(\mathbf{M}\) into \(\mathbf{N}\) by
\end{enumerate}

\[
h \left(x _ {1}, x _ {2}, \dots , x _ {q}\right) = \frac {1}{q !} \sum_ {\sigma \in \mathfrak {S} _ {q}} g \left(x _ {\sigma (1)}, x _ {\sigma (2)}, \dots , x _ {\sigma (q)}\right).
\]

Then \(h\) is symmetric and \(f(x) = h(x, x, \ldots, x)\) for all \(x \in \mathbf{M}\) .

\begin{enumerate}
\def\labelenumi{\alph{enumi})}
\setcounter{enumi}{1}
\tightlist
\item
  Let \(h\) be a symmetric \(q\) -linear mapping of \(\mathbf{M}^q\) into \(\mathbf{N}\) such that \(f(x) = g(x, x, \ldots, x)\) . Let \(l\) be the linear mapping of \(\mathbf{T}^q(\mathbf{M})\) into \(\mathbf{N}\) such that \(h(x_1, \ldots, x_q) = l(x_1 \otimes \ldots \otimes x_q)\) for any \(x_1, \ldots, x_q \in \mathbf{M}\) . We have
\end{enumerate}

\[
\begin{array}{l} (- 1) ^ {q} q! h (x _ {1}, \dots , x _ {q}) = (- 1) ^ {q} \sum_ {\sigma \in \mathfrak {S} _ {q}} h (x _ {\sigma (1)}, \dots , x _ {\sigma (q)}) = \\ = (- 1) ^ {q} l (\mathbf {s} (x _ {1} \otimes \dots \otimes x _ {q})) = \sum_ {\mathrm{H} \subset \{1, \dots , q \}} (- 1) ^ {\text {Card H}} l \left(\gamma_ {q} \left(\sum_ {i \in \mathrm{H}} x _ {i}\right)\right) \end{array}
\]

by Prop. 3, (v) (IV, p.~45). Now

\[
l \left(\gamma_ {q} \left(\sum_ {i \in H} x _ {i}\right)\right) = h \left(\sum_ {i \in H} x _ {i}, \dots , \sum_ {i \in H} x _ {i}\right) = f \left(\sum_ {i \in H} x _ {i}\right)
\]

and this proves formula (16) and the uniqueness of h.

PROPOSITION 16. --- Let M be a free A-module, N an A-module, q a positive integer and u the canonical homomorphism of Hom(TS \(^{q}\) (M), N) into Pol \(^{q}\) (M, N).

\begin{enumerate}
\def\labelenumi{(\roman{enumi})}
\item
  If \(\mathbf{A}\) is an infinite integral domain and \(\mathbf{N}\) is torsion-free, then \(u\) is an isomorphism.
\item
  If the mapping \(y \mapsto q!y\) in \(\mathbf{N}\) is injective, \(u\) is an isomorphism.
\end{enumerate}

In the two cases of the proposition we have to prove that \(u\) is injective, that is, every linear mapping \(f\) of \(\mathbf{TS}^q(\mathbf{M})\) into \(\mathbf{N}\) which is zero on \(\gamma_q(\mathbf{M})\) , vanishes.

Assume A to be an infinite integral domain and N torsion-free. For every \(z \in \mathbf{TS}^{q}(\mathbf{M})\) there exists \(\alpha \in A - \{0\}\) such that \(\alpha z\) is an A-linear combination of elements of \(\gamma_{q}(\mathbf{M})\) (IV, p.~47, Prop. 5). Hence \(\alpha f(z) = f(\alpha z) = 0\) , and so \(f(z) = 0\) .

Suppose next that the mapping \(y \mapsto q!y\) in \(\mathbf{N}\) is injective, then by IV, p.~45, Prop. 3, (v), \(f\) vanishes on \(\mathbf{s} \cdot \mathbf{T}^q(\mathbf{M})\) . Hence if \(z \in \mathbf{T}\mathbf{S}^q(\mathbf{M})\) , we have \(q!f(z) = f(\mathbf{s}z) = 0\) , and so \(f(z) = 0\) .

COROLLARY. --- Let M be a free A-module, N an A-module, q a positive integer, \(h \in \operatorname{Pol}^{q}(M, N)\) and \((e_{i})_{i \in I}\) a basis of M. In the two cases of Prop. 16 there exists a unique family \((u_{\nu})_{\nu \in \mathbf{N}^{(I)}, |\nu| = q}\) of elements of N such that \(h\left(\sum_{i \in I} \lambda_{i} e_{i}\right) = \sum_{|\nu| = q} \lambda^{\nu} u_{\nu}\) for all \((\lambda_{i}) \in \mathbf{A}^{(I)}\) .

